% Part: second-order-logic
% Chapter: sol-and-sets
% Section: comparing-sets

\documentclass[../../../include/open-logic-section]{subfiles}

\begin{document}

\olfileid{sol}{set}{cmp}

\olsection{संचांची तुलना}

\begin{prop}
!!{formula} $\lforall[x][(X(x) \lif Y(x))]$ उपसंच-संबंध
परिभाषित करते; म्हणजे
$\Sat{M}{\lforall[x][(X(x) \lif Y(x))]}[s]$
जर आणि फक्त जर $s(X) \subseteq s(Y)$.
\end{prop}

\begin{prop}
!!{formula} $\lforall[x][(X(x) \liff Y(x))]$ संचांवरील
समानता-संबंध परिभाषित करते; म्हणजे
$\Sat{M}{\lforall[x][(X(x) \liff Y(x))]}[s]$
जर आणि फक्त जर $s(X) = s(Y)$.
\end{prop}

\begin{prop}
!!{formula} $\lexists[x][X(x)]$ अरिक्त असण्याचा
गुणधर्म परिभाषित करते; म्हणजे
$\Sat{M}{\lexists[x][X(x)]}[s]$
जर आणि फक्त जर $s(X) \neq \emptyset$.
\end{prop}

संच~$X$ हा संच~$Y$ पेक्षा आकारमानाने मोठा नाही,
$\cardle{X}{Y}$, जर आणि फक्त जर
!!a{injective} फलन $f\colon X \to Y$ असेल.
फलन एकास-एक आहे, तसेच~$X$ मधील निविष्टांसाठी
त्याची मूल्ये~$Y$ मध्ये आहेत, हे व्यक्त करता येते.
म्हणून प्रांताच्या उपसंचांवर आकारमानाने मोठा नाही
हा संबंधही परिभाषित करता येतो.

\begin{prop}
पुढील सूत्र
\[
\lexists[u][(\lforall[x][(X(x) \lif Y(u(x)))] \land \lforall[x][\lforall[y][(\eq[u(x)][u(y)] \lif
      \eq[x][y])]])]
\]
आकारमानाने मोठा नाही हा संबंध परिभाषित करते.
\end{prop}

दोन संचांचे आकारमान समान असते, म्हणजे ते
``तुल्यबल'' असतात, $\cardeq{X}{Y}$, जर आणि
फक्त जर !!a{bijective} फलन~$f\colon X \to Y$ असेल.

\begin{prop}
पुढील सूत्र
\begin{multline*}
  \lexists[u][(\lforall[x][(X(x) \lif Y(u(x)))] \land {}\\
    \lforall[x][\lforall[y][(\eq[u(x)][u(y)] \lif \eq[x][y])]]
  \land {} \\\lforall[y][(Y(y) \lif \lexists[x][(X(x)
      \land \eq[y][u(x)])])])]
\end{multline*}
तुल्यबल असण्याचा संबंध परिभाषित करते.
\end{prop}

या !!{formula}s ना अनुक्रमे $X \subseteq Y$,
$X = Y$, $X \neq \emptyset$, $\cardle{X}{Y}$
आणि $\cardeq{X}{Y}$ अशी संक्षिप्त चिन्हे वापरू.
(संच $X$ आणि $Y$ यांच्याविषयी अनौपचारिकपणे
बोलतानाही हीच चिन्हे वापरल्यामुळे थोडा गोंधळ होऊ शकतो;
पण येथे ती एकस्थानी संबंधचले $X$ आणि~$Y$
असलेल्या द्वितीय-क्रम तर्कशास्त्रातील !!{formula}s
साठीची संक्षिप्त चिन्हे आहेत.)

\begin{prop}
!!{sentence} $\lforall[X][\lforall[Y][((\cardle{X}{Y} \land
    \cardle{Y}{X}) \lif \cardeq{X}{Y})]]$ वैध आहे.
\end{prop}

\begin{proof}
कोणत्याही उपसंचांसाठी $X \subseteq \Domain{M}$
आणि $Y \subseteq \Domain{M}$, जर
$\cardle{X}{Y}$ आणि $\cardle{Y}{X}$, तर
$\cardeq{X}{Y}$ असेल, तर !!a{structure}~$\Struct{M}$
मध्ये या !!{sentence} ची पूर्ति होते. पण हे
\emph{कोणत्याही} संच $X$ आणि $Y$ साठी खरे आहे:
हेच श्रेडर--बर्नस्टाइन प्रमेय आहे.
\end{proof}


\end{document}
